\subsection{Which test the gain comes from}\label{app:ablation}

\begin{table}[t]
\caption{Ablation of the controller's tests on the frozen checkpoint: phase-B throughput of each variant relative to \pol{Full} on the same request sequences (paired mean $\pm$ standard deviation over three repetitions).}
\label{tab:ablate}
\centering\small
\IfFileExists{tables/app_ablate.tex}{% generated by code/specsched/paper_assets.py -- do not edit by hand
\begin{tabular}{@{}lccc@{}}
\toprule
Variant & 8B code & 8B Chinese & MoE code\\
\midrule
Full without the payback rule & $-2.5$\,$\pm$2.6 & $+2.4$\,$\pm$4.2 & $+0.4$\,$\pm$1.7 \\
Full without the stop tests & $-0.2$\,$\pm$1.1 & $+1.8$\,$\pm$2.5 & $-0.9$\,$\pm$0.3 \\
Full without the speculation switch & $+0.2$\,$\pm$1.6 & $-43.8$\,$\pm$2.2 & $+0.1$\,$\pm$0.6 \\
Full without the per-domain knowledge & $-0.6$\,$\pm$2.7 & $-0.1$\,$\pm$6.6 & $-0.9$\,$\pm$1.5 \\
\bottomrule
\end{tabular}
}{}
\end{table}

The policies of \cref{tab:main} are partial ablations, but none isolates the payback rule, and \pol{Switch} removes training altogether. We therefore reran \pol{Full} with one of its tests removed, on the frozen checkpoint, prompts, and seeds (\nv{abl.n.runs} runs; Table~\ref{tab:ablate}): \emph{without the payback rule}, the trainer runs at duty 0.5 whenever the stop tests allow it under saturation; \emph{without the stop tests}, the structural-ceiling, at-ceiling, and convergence tests are skipped; \emph{without the switch}, speculation stays on and training is valued with speculation stuck on (\cref{app:controller}); \emph{without per-domain knowledge}, the controller ignores the traffic label and the learning curves: it predicts throughput with the per-target proportional frontier $\Ghat(a,C)=a\,T_{\mathrm{base}}(C)/T_{\mathrm{spec}}(C)$, assumes that the next round raises acceptance by 0.3, replaces that prior with the gain measured after every swap, and, after a round that did not pay, waits 120\,s before trying one more.

The gain comes from a different test in each scenario. On Chinese it comes entirely from the switch (\nv{abl.noswitch.zh.vsv6}\% without it, \nv{abl.noswitch.zh.min} to \nv{abl.noswitch.zh.max}\% per repetition), while the other three variants stay within noise of \pol{Full}. On Qwen3-8B code it comes from training, and the payback rule is worth \nv{abl.novalue.code.vsv6.abs}\% there: at a fixed duty the trainer runs until the ceiling test stops it (mean duty \nv{abl.novalue.code.duty} against \nv{code.full.duty}), and the interference costs more than the last drafters gain. On the MoE target every variant trains and stops at the ceiling, and no single test is decisive. The stop tests change nothing under saturation (\nv{abl.nostop.code.vsv6}\%, \nv{abl.nostop.zh.vsv6}\%, \nv{abl.nostop.moe.vsv6}\%): near the ceiling the payback rule already finds no value, and on Chinese it finds none even without the structural-ceiling test, so the three tests are shortcuts that matter only for the SLO-slack rule in the light phases. The variant without per-domain knowledge stays within noise of \pol{Full} in all three scenarios (\nv{abl.nocurve.code.vsv6}\%, \nv{abl.nocurve.zh.vsv6}\%, \nv{abl.nocurve.moe.vsv6}\%): the switch needs only the measured acceptance and the per-target step times, and the measured swap gains replace the learning curve. Its one cost is timing: in one Chinese repetition (\nv{abl.nocurve.zh.min}\%) a drafter trained in the slack reached its export point late in phase B, and the speculation probe after the swap ran at saturation on Chinese. Where the curves matter is traffic that a label does not describe, which \cref{app:mix} examines.
